\begin{recipe}
[% 
    preparationtime = {\unit[40]{min}},
    portion = {\portion{1}},
    source = {Die Anatomie des guten Kochens s.73 \& Eigenkreation },
    calory = {875kcal | 19g Protein | 92g KH | 42g Fett},
    price = {4,66€},
    rating = {5},
    created = {11.03.2026}
]
{Zitronen-Risotto mit gegrilltem Romanaherz}

\graph{% pictures
    big=pic/hauptgericht/10_zitronen-risotto-mit-gegrilltem-romana
}

\introduction{%
Das Rezept des gegrillte Romanaherzens\index{Gemüse!Romana} habe ich dem genannten Buch aus dem Kapitel Säure, deswegen dachte ich mir, dass zum satt machen, ein Zitronenrisotto ideal passt. Durch den Salat hat das Gericht noch eine weitere frische Komponente, die sehr gut zu dem Risotto harmoniert.
}

\ingredients{
    \textbf{Risotto} & \\
    1 & kleine Zwiebel\\
    100 g & Risottoreis\index{Getreide!Risottoreis}\\
    85 ml & Weißwein\\
    300 ml & Gemüsebrühe\\
    2 EL & Zitronesaft\index{Obst!Zitrone}\\
    2 EL & Hefeflocken\\
    \\
    \textbf{Kürbiskerne} & \\
    20 g & Kürbiskerne\index{Gemüse!Kürbis}\\
    \nicefrac{1}{2} EL & Olivenöl\\
    \nicefrac{1}{2} TL & Paprikapulver\\
    \nicefrac{1}{2} TL & Chilipluver\\
    \\
    \textbf{Creme} & \\
    60 g & Joghurt\index{Milchprodukte \& Alternativen!Joghurt}\\
    25 g & Crème fraîche\index{Milchprodukte \& Alternativen!Crème fraîche}\\
    1 EL & Granatapfelsirup\\
    \nicefrac{1}{2} TL & Kurkuma\\
    2 EL & Limettensaft\index{Obst!Limette}\\
    1 Zehe & Knoblauch\index{Gemüse!Knoblauch}\\
    \multicolumn{2}{l}{\textit{\small Fortsetzung auf nächster Seite}}
}
\ingredientsB{
    \textbf{Salat} & \\
    1 & Romana-Salatherzen\\
    & Olivenöl zum Grillen\\
    & Meersalz
}
\preparation{%
    \step%
    Zwiebel fein schneiden und in einem Topf mit Olivenöl 1--2min bei mittlerer Hitze anschwitzen
    
    \step%
    Den Risottoreis zugeben und weitere 2min anschwitzen. Mit Weißwein ablöschen und einkochen lassen.
    
    \step%
    Nach und nach heiße Gemüsebrühe zugeben und unter ständigem Rühren einkochen lassen, bis der Reis cremig und gar ist.
    
    \step%
    Zitronenabrieb und etwas Zitronensaft unterrühren und mit Salz und Pfeffer abschmecken.
    
    \step%
    Für die Creme Joghurt, Crème fraîche, Granatapfelsirup, Kurkuma, Limettensaft und Knoblauch glatt rühren.
    
    \step%
    Romana-Salatherzen halbieren, mit Olivenöl bestreichen und in einer heißen Grillpfanne mit der Schnittfläche nach unten grillen, bis deutliche Röstaromen entstehen.
    
    \step%
    Risotto auf Tellern anrichten, gegrillten Romana darauf platzieren mit der Creme beträufeln und die Kürbiskerne dazu streuen.
}

\hint{%
    Bestenfalls einen dunklen Teller verwenden, falls ein heller verwendet wird, kann man auch noch mit Balsamico Tropfen hinzugeben. Falls gewünscht auch auf den Salat, aber nur in geringen Mengen, das es sonst zu säuerlich wird und die Zitrone nicht mehr herauskommt.
}

\suggestion[Menge Romanaherz]{%
    Auf dem Bild sieht man nur eine Häfte des Salates, da sonst das Risotto nicht mehr gut zu sehen ist. Wenn es ein ansehnliches Gericht werden soll, dann empfehle ich auch nur die Hälfte zu nehmen. Ein ganzes Romanaherz sollte nur bei großem Hunger und den 100g Reis genommen werden, da es sonst sich nicht gleichmäßig essen lässt und es wirkt einfach nach zu viel Salat.
}

\end{recipe}